\documentclass[10pt,a4paper]{article}
\usepackage[T1]{fontenc}
\usepackage[utf8]{inputenc}
\usepackage{lmodern}
\usepackage[spanish]{babel}
\usepackage[margin=1.7cm]{geometry}
\usepackage{tikz}
\usetikzlibrary{arrows.meta,positioning}
\usepackage{hyperref,fancyhdr,booktabs,tabularx,listings}
\hypersetup{hidelinks}
\pagestyle{fancy}\fancyhf{}
\fancyhead[L]{Vía · EI1039 / EI1048}
\fancyhead[R]{Entrega 2 · Pruebas y SUT}
\fancyfoot[C]{\thepage}
\setlength{\headheight}{14pt}
\setlength{\parindent}{0pt}
\setlength{\parskip}{6pt}
\lstset{basicstyle=\ttfamily\fontsize{8}{9.5}\selectfont,breaklines=true,columns=fullflexible,keepspaces=true,showstringspaces=false,frame=single,rulecolor=\color{black!25},backgroundcolor=\color{black!2},aboveskip=5pt,belowskip=8pt}
\tikzset{box/.style={draw,align=center,text width=6.0cm,font=\small,inner sep=5pt,fill=blue!3},dep/.style={dashed,-{Stealth[length=2mm]}},real/.style={dashed,-{Triangle[open,length=2.5mm]}},lab/.style={font=\scriptsize,fill=white,inner sep=2pt}}
\newcommand{\codigo}[2]{\href{https://github.com/Kiliansito/ei1039-ei1048-movilidad/blob/7c37466/#1}{\texttt{\small #2}}}
\begin{document}
{\LARGE\bfseries Iteración 1: pruebas y SUT}\par
{\large Apartados 2 y 3 de la Entrega 2}\par

\section*{Alcance y estado}
La propuesta de IT1 comprende HU-01 (registro), HU-02 (inicio de sesión), HU-05 (alta de lugares por coordenadas), HU-07 (consulta de lugares) y HU-14 (cálculo de rutas). Su distribución definitiva en la planificación global se acordará con el equipo.

Se han programado 25 pruebas para los escenarios identificados en la Entrega 1. El SUT mínimo define constructores, métodos, tipos y dependencias para poder ejecutarlas. Los métodos nuevos lanzan un error de funcionalidad pendiente. Esta ejecución documenta la fase roja de ATDD; la implementación funcional y la fase verde corresponden al desarrollo posterior.

\begin{center}
\begin{tabularx}{\linewidth}{l c X}
\toprule
Historia & Pruebas & Escenarios identificados\\
\midrule
HU-01 & 6 & Un registro válido y cinco escenarios inválidos.\\
HU-02 & 5 & Un acceso válido y cuatro escenarios inválidos.\\
HU-05 & 5 & Un alta válida y cuatro escenarios inválidos.\\
HU-07 & 4 & Dos consultas válidas y dos escenarios inválidos.\\
HU-14 & 5 & Un trayecto válido y cuatro escenarios inválidos.\\
\midrule
Total & 25 & Resultado de ejecución: 0 pasan, 25 fallan.\\
\bottomrule
\end{tabularx}
\end{center}

\section*{Ejecución reproducible y evidencia}
Desde la raíz del repositorio, con Node.js compatible con los módulos y la transformación TypeScript usados por el spike, ejecutar en PowerShell:
\begin{lstlisting}
$pruebas = Get-ChildItem tests -Filter '*.aceptacion.test.ts' | Select-Object -ExpandProperty FullName
node --experimental-transform-types --env-file-if-exists=entrega1/02_Spike/.env --test-reporter=spec --test $pruebas
\end{lstlisting}

La evidencia conserva versión de Node, fecha UTC, commit de código y salida completa. Se verificó que las 25 pruebas se ejecutan, que el proceso devuelve código 1 y que los fallos proceden de los métodos pendientes, sin errores de carga o sintaxis. \href{https://github.com/Kiliansito/ei1039-ei1048-movilidad/blob/h-14/docs/entrega2/evidencias/pruebas-it01-rojo.txt}{Ver evidencia de ejecución}.

\textbf{Dependencias.} Las pruebas de datos usan SQLite real en memoria y preparan sus propios datos. El caso válido de ruta recibe el adaptador real de ORS y necesitará \texttt{ORS\_API\_KEY} al implementar el servicio. En la fase roja todavía no llega a consultar ORS.

\textbf{Indisponibilidad.} Los casos de conexión preparan un servidor HTTP local que corta la conexión. En E-14-I3 se redirige la dependencia HTTP del adaptador a ese servidor; se trata de una comprobación de integración con el fallo controlado de un sustituto del proveedor. Esta clasificación debe conservarse en la documentación de pruebas.

\textbf{Código completo.}
\codigo{tests/registro.aceptacion.test.ts}{HU-01: registro};
\codigo{tests/inicio-sesion.aceptacion.test.ts}{HU-02: sesion};
\codigo{tests/guardar-lugar.aceptacion.test.ts}{HU-05: alta de lugar};
\codigo{tests/consultar-lugares.aceptacion.test.ts}{HU-07: consulta de lugares};
\codigo{tests/calcular-ruta.aceptacion.test.ts}{HU-14: rutas}.

\newpage
{\LARGE\bfseries Ejemplos de pruebas de aceptación}\par
Los ejemplos siguientes reproducen dos pruebas de la HU-01. Antes de cada prueba se crea una base SQLite vacía con la tabla \texttt{cuentas}. \texttt{datosDeAna} contiene nombre Ana, correo \texttt{ana@example.com} y contraseña \texttt{RutaSegura2026}. Los imports, el esquema y la preparación completa están en el archivo enlazado.

\section*{PA-01-V: registro válido}
\lstinputlisting[firstline=34,lastline=48]{../../tests/registro.aceptacion.test.ts}

Se comprueban la creación de una única cuenta, su correo y que la respuesta contiene únicamente el identificador de la cuenta guardada.

\section*{PA-01-I: correo duplicado}
\lstinputlisting[firstline=50,lastline=68]{../../tests/registro.aceptacion.test.ts}

Se prepara una cuenta existente, se intenta repetir el alta y se exige un error de correo duplicado. La comprobación final conserva una sola cuenta con su identificador original. Cada prueba prepara su estado; la prueba inválida utiliza una inserción directa para evitar depender de que funcione el registro.

Los nombres siguen \texttt{metodoBajoPruebas\_escenario\_comportamientoEsperado}. Los comentarios separan Dado, Cuando y Entonces. En la fase actual ambas pruebas fallan al alcanzar \texttt{registrar}, que lanza \texttt{Registro sin implementar}.

\newpage
{\LARGE\bfseries UML del sistema bajo prueba}\par
Las clases se corresponden con \codigo{src/ServicioCuentas.ts}{ServicioCuentas},
\codigo{src/ServicioLugares.ts}{ServicioLugares},
\codigo{src/ServicioRutas.ts}{ServicioRutas},
\codigo{src/ClienteCuentas.ts}{ClienteCuentas} y
\codigo{src/ClienteLugares.ts}{ClienteLugares}.
Las flechas discontinuas representan dependencias; la punta triangular vacía indica implementación de la interfaz.

\begin{center}
\begin{tikzpicture}
\node[box] (cc) at (0,0) {\textbf{ClienteCuentas}\\constructor(urlBase)\\registrar(datos):\\Promise de ResultadoRegistro\\iniciarSesion(credenciales):\\Promise de ResultadoSesion};
\node[box] (cl) at (8,0) {\textbf{ClienteLugares}\\constructor(urlBase)\\consultarLugares(token): Promise de Lugar[]};
\node[box] (sc) at (0,-3.3) {\textbf{ServicioCuentas}\\constructor(baseDeDatos)\\registrar(datos): ResultadoRegistro\\iniciarSesion(credenciales): ResultadoSesion};
\node[box] (sl) at (8,-3.3) {\textbf{ServicioLugares}\\constructor(baseDeDatos)\\guardarPorCoordenadas(token, datos): Lugar\\consultarLugares(token): Lugar[]};
\node[box] (sr) at (0,-6.6) {\textbf{ServicioRutas}\\constructor(baseDeDatos, proveedor)\\calcularRuta(token, datos):\\Promise de RutaCalculada};
\node[box] (db) at (8,-6.6) {$\ll$SQLite$\gg$\\\textbf{DatabaseSync}};
\node[box] (rp) at (0,-9.2) {$\ll$interface$\gg$\\\textbf{RoutingPort}\\routes(from, to, mode):\\Promise de RawRoute[]};
\node[box] (ors) at (8,-9.2) {\textbf{OrsAdapter}\\routes(from, to, mode)};
\draw[dep] (cc) -- node[lab,right]{tipos} (sc);
\draw[dep] (cl) -- node[lab,right]{tipo Lugar} (sl);
\draw[dep] (sc.south) -- ++(0,-.65) -| (db.west);
\draw[dep] (sl) -- (db);
\draw[dep] (sr) -- (db);
\draw[dep] (sr) -- (rp);
\draw[real] (ors) -- (rp);
\end{tikzpicture}
\end{center}

\textbf{Datos y resultados.}
\texttt{DatosRegistro} agrupa nombre, correo y contraseña.
\texttt{ResultadoRegistro} devuelve el identificador.
\texttt{Credenciales} agrupa correo y contraseña.
\texttt{ResultadoSesion} devuelve token y cuenta pública.
\texttt{Lugar} contiene identificador, nombre, latitud y longitud.
\texttt{SolicitudRuta} contiene los identificadores de origen y destino y el modo.
\texttt{RutaCalculada} contiene geometría, distancia en metros y duración en segundos.
Estos nombres abrevian los tipos definidos en línea en el código.

Los constructores reciben las dependencias que preparan las pruebas. Las dependencias de los clientes a los servicios son referencias de tipos. El adaptador y su contrato se reutilizan del spike. Los métodos del SUT nuevo permanecen pendientes de implementación; el diagrama muestra sus firmas y las dependencias declaradas.
\end{document}
